\section{Detailed proofs}
\label{app:proofs}

This appendix collects the detailed proofs of lemmas whose
inline statements in \S\ref{sec:lemmas} are accompanied by
proof outlines.  The appendix proofs are self-contained.

\subsection{Proof of Lemma 1 (Intensive composition under co-evolution)}
\label{app:proof_lem1}

\textbf{Statement (restated, generic form).}  Let $X$ be a non-residual
gap quantity (e.g.,\ $\epsnonresComposed$, the composed
non-residual gap under co-evolution; or $f(\epsilon_{\mathrm{safe}})$
after Lemma~\ref{lem:2}'s substitution).  The composition rule
\[
  B \;=\; \mathrm{Lip}(g) \cdot \bigl(X + \lambda \cdot \epsfloor\bigr)
\]
where $\lambda \in [0, 1]$ is the residual weight and
$\mathrm{Lip}(g)$ is the Lipschitz constant of the alignment
property $g$, is intensive in $|P|$ over a deployment class
$\mathcal{D}$ provided $X$, $\epsfloor$, and $\mathrm{Lip}(g)$ are
each intensive in $|P|$ over $\mathcal{D}$, and the co-evolution
correction terms in the structure of $X$ are bounded by Paper~10's
Assumption~\ref{ass:bounded_coev} (bounded co-evolution).

\begin{proof}
We work over a deployment class $\mathcal{D}$ characterized by
fixed operational parameters (threat model, training-discipline
regime, verification infrastructure, governance protocol).  The
operational parameters are fixed before $|P|$ is varied; they may
not include arbitrary state-dependent quantities chosen to absorb
the bound.

\emph{Step 1 (per-capability admissibility).}  Under
condition~(C8) of Theorem~\ref{thm:deployment_safety} (per-capability
admissibility, inherited from
\citep[Microfoundation]{lasser2026micro}'s Assumption~1), for every
capability $c \in P^{\mathrm{act}}$ in the operationally active
subspace, the per-capability proxy-truth gap is allocated a ceiling
$\kappa_c$, and each channel $j \in \{1, 2, 3, 4\}$ receives a
disjoint sub-share $\kappa_c^{(j)}$ of the ceiling, with
$\sum_j \kappa_c^{(j)} \leq \kappa_c$.

\emph{Step 2 (per-channel sup-norm bridge).}  Define the per-channel
non-residual gap as the sup-norm of the channel-$j$ contributions
across capabilities:
\[
  \epsnonresChannel{j} \;:=\; \sup_{c \in P^{\mathrm{act}}}
  \bigl(\text{channel-}j\text{ contribution to } |T(c) - P(c)|\bigr).
\]
By the disjoint sub-share allocation,
\[
  \epsnonresChannel{j} \;\leq\; \sup_{c \in P^{\mathrm{act}}}
  \kappa_c^{(j)} \;=:\; K_j,
\]
where $K_j$ depends on $\mathcal{D}$'s operational parameters
(specifically, the verification-infrastructure thresholds that set
the channel sub-share ceilings) but is independent of $|P|$. The
disjoint structure prevents double-counting across channels.

\emph{Step 3 (composed non-residual gap under co-evolution).}
Under the co-evolution regime, \citep[Microfoundation]{lasser2026micro}'s
Composition Proposition (in the lineage section) establishes
exact composition of the four channels in the sequential-intervention
regime, with simultaneous co-evolution introducing positive-part
correction terms:
\[
  \epsnonresComposed \;\leq\; \sum_{j \in \{1, 2, 3, 4\}}
  \epsnonresChannel{j} + (\epsnonresCoev)_+.
\]
The source paper leaves the formal characterization of
$(\epsnonresCoev)_+$ as open work; Paper~10 closes the gap with
Assumption~\ref{ass:bounded_coev} (bounded co-evolution).

\emph{Step 4 (bounded co-evolution).}  By
Assumption~\ref{ass:bounded_coev}, there exists
$\bar{M} \geq 0$ depending on $\mathcal{D}$'s operational parameters
but independent of $|P|$, such that the maximum per-channel
coupling magnitude satisfies $M(\text{deployment}) \leq \bar{M}$
for all valid deployment states.  Combined with the structural
co-evolution constant $K_{\mathrm{coev}} \geq 0$:
\[
  (\epsnonresCoev)_+ \;\leq\; K_{\mathrm{coev}} \cdot \bar{M}
  \;=:\; K'_{\mathrm{coev}}.
\]
$K'_{\mathrm{coev}}$ is independent of $|P|$.

\emph{Step 5 (bounded composed non-residual gap).}  Combining
Steps 2--4:
\[
  \epsnonresComposed \;\leq\; \sum_j K_j + K'_{\mathrm{coev}}
  \;=:\; K_{\mathrm{nonres}}.
\]
For the generic form, set $X \leq K_X$ where $K_X$ is the corresponding
$|P|$-independent constant (e.g., $K_X = K_{\mathrm{nonres}}$ for
$X = \epsnonresComposed$, or $K_X = f(\epsilon_{\mathrm{safe}})$ for
$X = f(\epsilon_{\mathrm{safe}})$ after Lemma~\ref{lem:2}).

\emph{Step 6 (residual floor).}  By the gap-decomposition structure
of \citep[Microfoundation]{lasser2026micro}, the residual floor
satisfies
$\epsfloor = \volR(\ResS)/\volR(\Pproxy) \in [0, 1]$
under finite positive $\volR(\Pproxy)$, with
$K_{\mathrm{floor}} \leq 1$ uniformly.

\emph{Step 7 (composed slack term).}  Define the composed slack term
\[
  X + \lambda \cdot \epsfloor \;\leq\; K_X + \lambda \cdot K_{\mathrm{floor}}
  \;\leq\; K_X + K_{\mathrm{floor}} \;=:\; K_{\mathrm{slack}}.
\]
$K_{\mathrm{slack}}$ is independent of $|P|$.  This is the central
result: the composed slack term itself is intensive.

\emph{Step 8 (Lipschitz multiplication).}  Under condition~(C6) of
Theorem~\ref{thm:deployment_safety} (bounded-Lipschitz alignment
property), $\mathrm{Lip}(g) \leq K_{\mathrm{Lip}}$ over $\mathcal{D}$.
Therefore:
\[
  B \;=\; \mathrm{Lip}(g) \cdot (X + \lambda \cdot \epsfloor)
  \;\leq\; K_{\mathrm{Lip}} \cdot K_{\mathrm{slack}} \;=:\; C^*.
\]
$C^*$ is the product of $|P|$-independent constants and is therefore
$|P|$-independent over $\mathcal{D}$.

By Definition~\ref{def:intensive_paper10} (intensive in $|P|$ over
$\mathcal{D}$), $B$ is intensive over $\mathcal{D}$. $\Box$
\end{proof}

\paragraph{Discussion of constants.}
The proof produces $C^* = K_{\mathrm{Lip}} \cdot (K_X +
K_{\mathrm{floor}})$, where for $X = \epsnonresComposed$ we have
$K_X = \sum_j K_j + K'_{\mathrm{coev}}$.  Each constant has an
operational interpretation:
\begin{itemize}
\item $K_{\mathrm{Lip}}$: bounded Lipschitz constant of $g$
  (\emph{deployment-policy-derived}; condition (C6) of
  Theorem~\ref{thm:deployment_safety}).
\item $K_j$ for $j \in \{1,2,3,4\}$: per-channel admissibility
  ceilings (\emph{source/admissibility-derived}, conditional on
  per-capability admissibility (C8) and channel sub-share
  allocation).
\item $K_{\mathrm{coev}}$: structural constant for the co-evolution
  dynamics (\emph{newly assumed}; not derived from a source-paper
  proposition).
\item $\bar{M}$: uniform bound on per-channel coupling magnitudes
  (\emph{newly assumed}; condition (C7) of
  Theorem~\ref{thm:deployment_safety}, the bounded co-evolution
  assumption).
\item $K_{\mathrm{floor}} \leq 1$: residual share bound
  (\emph{source-derived by ratio}, assuming finite positive
  $\volR(\Pproxy)$).
\end{itemize}

The deployment-tooling specification of \S\ref{sec:deployment_tooling}
must include calibration hooks for the newly-assumed constants
$K_{\mathrm{coev}}$ and $\bar{M}$ before claiming the deployment
class verifies (C7).  The other constants inherit calibration from
the source-paper machinery and the deployment's policy choices.

\subsection{Proof of Lemma 2 (Lyapunov-Goodhart bridge)}
\label{app:proof_lem2}

\textbf{Statement (restated).}  Under conditions (C9.BD) bounded
per-capability dimension count, (C9.CL) coordinate-Lipschitz
parameterization, (C9.WB) safety-relevant subspace with weight
floor, and (C9.TF) truth floor of
Theorem~\ref{thm:deployment_safety}, and under invariants $I_1$ and
$I_4$ (which together imply $I_2$):
\[
  \Lyap(\Wm_t) \;<\; \epsilon_{\mathrm{safe}}
  \quad \implies \quad
  \epsnonres \;<\; f(\epsilon_{\mathrm{safe}})
  \;:=\; \frac{L_{\max}}{T_{\min}}
  \sqrt{\frac{N_{\max} \cdot \epsilon_{\mathrm{safe}}}{w_{\min}}},
\]
where $\epsnonres = \sup_{c \in P^{\mathrm{act}}} |P(c) - T(c)|/T(c)$
is \citep[Microfoundation]{lasser2026micro}'s relative sup-norm gap
on the operationally active subspace
$P^{\mathrm{act}} = \{c \in P \setminus \mathrm{ResS} : T(c) > 0\}$.

\begin{proof}
We work over the deployment class $\mathcal{D}$ satisfying
conditions (C9.BD), (C9.CL), (C9.WB), and (C9.TF) throughout.

\emph{Step 1 (per-capability absolute gap).}  By (C9.CL), for
every capability $c \in P^{\mathrm{act}}$:
\[
  |P(c) - T(c)| \leq \sum_{k \in \mathrm{dim}(c)} L_{c,k} \cdot
  |\epsilon_k|.
\]

\emph{Step 2 (weighted-dual-norm Cauchy-Schwarz).}  Apply
Cauchy-Schwarz with weights $a_k = L_{c,k}/\sqrt{w_k}$ and
$b_k = \sqrt{w_k} |\epsilon_k|$ for $k \in \mathrm{dim}(c)$:
\[
  \sum_{k \in \mathrm{dim}(c)} L_{c,k} |\epsilon_k|
  \;\leq\;
  \sqrt{\sum_{k \in \mathrm{dim}(c)} \frac{L_{c,k}^2}{w_k}}
  \cdot \sqrt{\sum_{k \in \mathrm{dim}(c)} w_k \epsilon_k^2}.
\]
Since $\mathrm{dim}(c) \subseteq S$ (the safety-relevant coordinate
set) and the Lyapunov function $\Lyap = \sum_{k \in S} w_k
\epsilon_k^2$ extends over $S$ with strictly positive weights
(C9.WB inherited from \citep[Phase~Redundancy]{lasser2026phase}),
the second factor is bounded by $\sqrt{\Lyap}$:
\[
  \sum_{k \in \mathrm{dim}(c)} w_k \epsilon_k^2
  \;\leq\; \sum_{k \in S} w_k \epsilon_k^2
  \;=\; \Lyap.
\]

\emph{Step 3 (closed-form via $L_{\max}, w_{\min}, N_{\max}$).}
For the appendix-friendly closed form, bound the dual-norm factor
by deployment-class-uniform constants:
\[
  \sum_{k \in \mathrm{dim}(c)} \frac{L_{c,k}^2}{w_k}
  \;\leq\; \frac{L_{\max}^2}{w_{\min}} \cdot |\mathrm{dim}(c)|
  \;\leq\; \frac{L_{\max}^2 \cdot N_{\max}}{w_{\min}},
\]
using $L_{c,k} \leq L_{\max}$, $w_k \geq w_{\min}$ for $k \in S$
(both deployment-class constants), and $|\mathrm{dim}(c)| \leq
N_{\max}$ by (C9.BD).

\emph{Step 4 (sup over $P^{\mathrm{act}}$).}  Combining Steps 1--3:
\[
  \sup_{c \in P^{\mathrm{act}}} |P(c) - T(c)|
  \;\leq\; L_{\max} \sqrt{\frac{N_{\max} \cdot \Lyap}{w_{\min}}}.
\]

\emph{Step 5 (apply truth floor for relative sup-norm).}  By
(C9.TF), $T(c) \geq T_{\min} > 0$ for all $c \in P^{\mathrm{act}}$.
Therefore $|P(c) - T(c)|/T(c) \leq |P(c) - T(c)|/T_{\min}$, and:
\[
  \epsnonres \;=\; \sup_{c \in P^{\mathrm{act}}}
  \frac{|P(c) - T(c)|}{T(c)}
  \;\leq\; \frac{L_{\max}}{T_{\min}}
  \sqrt{\frac{N_{\max} \cdot \Lyap}{w_{\min}}}.
\]

\emph{Step 6 (apply Lyapunov bound).}  Under invariants $I_1, I_4$,
\citep[Phase~Redundancy]{lasser2026phase}'s contraction analysis
gives $\Lyap(\Wm_t) < \epsilon_{\mathrm{safe}}$ in the
self-correcting basin.  Substituting:
\[
  \epsnonres \;<\; \frac{L_{\max}}{T_{\min}}
  \sqrt{\frac{N_{\max} \cdot \epsilon_{\mathrm{safe}}}{w_{\min}}}
  \;=\; f(\epsilon_{\mathrm{safe}}). \qquad\Box
\]
\end{proof}

\paragraph{Tighter heterogeneous form.}
The closed-form bound uses uniform constants $L_{\max}, w_{\min}$;
the underlying weighted-dual-norm bound (Step~2) is sharper when
deployment-specific sensitivities are heterogeneous:
\[
  \epsnonres \;\leq\;
  \sup_{c \in P^{\mathrm{act}}}
  \frac{1}{T(c)}
  \sqrt{\sum_{k \in \mathrm{dim}(c)} \frac{L_{c,k}^2}{w_k}}
  \cdot \sqrt{\Lyap}.
\]
With a finite $T_{\max} = \sup_{c} T(c)$, the looseness of the
closed form is bounded by approximately $T_{\max}/T_{\min}$;
without a truth ceiling, the looseness is not uniformly bounded
but the closed form remains a correct upper bound.  Operators may
use the heterogeneous form for deployment-specific calibration.

\paragraph{Discussion of constants.}
\begin{itemize}
\item $\epsilon_{\mathrm{safe}}$:
  \emph{source-derived} from
  \citep[Phase~Redundancy]{lasser2026phase}'s policy-correctness
  threshold.
\item $w_{\min}$: \emph{source-inherited and deployment-policy-set}.
  PR's Lyapunov weights are $w_k > 0$ on the safety-relevant
  subspace $S$; the floor $w_{\min} = \min_{k \in S} w_k > 0$ is
  set by the deployment's safety-relevance weighting.
\item $N_{\max}$: \emph{deployment-policy-derived} (C9.BD).
\item $L_{\max}, L_{c,k}$: \emph{deployment-policy-derived} (C9.CL).
\item $T_{\min}$: \emph{deployment-policy-derived and
  deployment-verified} (C9.TF).
\end{itemize}

\subsection{Proof of Lemma 5a (Substrate floor)}
\label{app:proof_lem5a}

\textbf{Statement (restated).}  Under invariant $I_6'$ ($\mindep
\geq m^* \geq 3$), condition (C10) of
Theorem~\ref{thm:deployment_safety} (per-pair cooperative-novelty
rate floor with superposition; comprising sub-clauses (C10.CN)
heterogeneous per-pair rate floors $\rho_{ij} \geq \rho_0 > 0$ on
audited subset $\Saudit$ of $m^*$ substrates, and (C10.SU)
pairwise channel superposition: disjoint attribution +
non-rivalrous production + joint-deployment intensities):
\[
  \rext \;\geq\; r_*(m^*) \;:=\; \rho_0 \cdot \binom{m^*}{2} \;>\; 0.
\]

\paragraph{Pairwise cooperative channels.}
For audited subset $\Saudit \subseteq \mathcal{S}$ with $|\Saudit|
= m^*$, the \emph{pairwise cooperative channel}
$\mathcal{C}^{(s_i, s_j)}$ for $i < j$ contains cooperatives with
exact two-substrate causally-necessary participation: cooperatives
whose production requires participants from $s_i$ and $s_j$ but
not from any other substrate.  Auxiliary participants whose role
is post-production (e.g., a Formal-Operational verifier attesting
to a Human-AI cooperative output) do not change the channel
classification; auxiliary participants whose role is causally
necessary for production (e.g., Formal-Operational rule-execution
that the output depends on) move the cooperative to the
appropriate higher-order channel.

\paragraph{Standing measure conventions.}
Throughout this proof:
\begin{itemize}
\item $\Saudit$ is an audit-time-certified subset, fixed for the
  deployment window/epoch.  Re-auditing across epochs is allowed
  per Audit~6 of \S\ref{sec:tooling_substrate_audits}; the bound
  applies per epoch or as the infimum across the deployment window.
\item Pairwise cooperative event classes are measurable subsets of
  $P^{\mathrm{act}}$.  Per-pair rates $\rho_{ij}$ are measured over
  a common time/normalization window matching $\rext$'s definition
  in \citep[Horizon~Aware]{lasser2026horizon}.
\item $\rext$ is treated as an additive nonnegative
  rate/measure over disjoint event classes.
\item Higher-order cooperative contributions are nonnegative.
\end{itemize}

\begin{proof}
We work over deployment class $\mathcal{D}$ satisfying $I_6'$ and
(C10).

\emph{Step 1 (substrate audit).}  By $I_6'$, $\mindep \geq m^*$
substrates with joint failure-correlation independence.  By
(C10), audited subset $\Saudit$ with $|\Saudit| = m^*$ is
certified for the deployment window.

\emph{Step 2 (combinatorial pairwise channels).}  The number of
unordered pairs in $\Saudit$ is $\binom{m^*}{2}$.  Each pair
$(s_i, s_j)$ with $i < j$ has channel $\mathcal{C}^{(s_i, s_j)}$
with exact two-substrate support.

\emph{Step 3 (per-pair rate from C10.CN).}  Each pair satisfies
$\rext^{(s_i, s_j)} \geq \rho_{ij} \geq \rho_0 > 0$ where $\rho_0
= \min_{i<j \in \Saudit} \rho_{ij}$.  By the joint-deployment-
intensity sub-clause (within (C10.CN)'s calibration clause and
(C10.SU)(3)), $\rho_{ij}$ measures the per-pair rate as realized
in the actual deployment, not under counterfactual isolation.

\emph{Step 4 (additivity from C10.SU).}  By (C10.SU)(1) (disjoint
attribution; preserved by exact-two-substrate support), the
pairwise channels are pairwise disjoint event classes.  By
(C10.SU)(2) (non-rivalrous production), simultaneous production
across pairs does not reduce any individual rate.  Therefore the
pairwise rates sum:
\[
  \rext^{\mathrm{pairwise}} \;:=\; \sum_{i<j \in \Saudit}
  \rext^{(s_i, s_j)} \;\geq\; \binom{m^*}{2} \cdot \rho_0.
\]

\emph{Step 5 (lower bound on $\rext$).}  Total $\rext$ is the
nonnegative additive measure over all cooperative event classes,
including pairwise plus higher-order.  Higher-order contributions
are nonnegative by convention, so:
\[
  \rext \;\geq\; \rext^{\mathrm{pairwise}} \;\geq\; \binom{m^*}{2}
  \rho_0 \;=\; r_*(m^*).
\]

\emph{Step 6 (intensivity over $\mathcal{D}$).}  $\rho_0$ is a
deployment-class constant (C10.CN, calibrated per epoch).
$m^*$ is the $I_6'$ threshold, deployment-class.  Therefore
$r_*(m^*)$ is independent of $|P|$ over $\mathcal{D}$.
$\Box$
\end{proof}

\paragraph{Discussion of constants.}
\begin{itemize}
\item $m^*$: \emph{deployment-policy-derived} threshold from $I_6'$.
\item $\rho_{ij}$: \emph{per-pair calibrated rates} (C10.CN);
  measured in joint deployment via Audit~6.
\item $\rho_0 = \min_{i<j \in \Saudit} \rho_{ij}$: derived from
  $\rho_{ij}$ via the audited-subset minimum.
\item $\Saudit$: \emph{audit-time certified subset}; may be
  re-audited across epochs.
\end{itemize}

\paragraph{On scope and tightness.}
The pairwise-only bound is conservative: higher-order cooperatives
contribute additionally to $\rext$.  Operators may calibrate
higher-order rates $\rho_0^{(k)}$ for tighter bounds, giving
$\rext \geq \sum_k \rho_0^{(k)} \binom{m^*}{k}$.  The lemma's
pairwise floor suffices to make Layer~1's safe region non-trivial
under any deployment-class $\rho_0 > 0$.

\paragraph{On (C10.SU) vs.\ (C2) joint failure-correlation independence.}
$I_6'$'s joint failure-correlation independence excludes shared
adversarial \emph{shock} mechanisms: a single shock cannot span
multiple substrates.  (C10.SU) excludes shared cooperative
\emph{production} bottlenecks: producing cooperatives in one
pairwise channel does not consume resources that would otherwise
contribute to another.  These are structurally independent
conditions; (C10.SU) is new content beyond $I_6'$, surfaced as a
deployment-class assumption.

\subsection{Proof of Lemma 4 (SPRT lead-time tail bound)}
\label{app:proof_lem4}

\textbf{Statement (restated).}  Let $T_{\mathrm{detect}}$ be the
SPRT detection time for a violation producing behavioral KL divergence
$\delta > 0$ under the alternative distribution. Let
$T_{\mathrm{cascade}}$ be the time from violation onset to
irreversible cascade. Then:
\[
  \Pr[T_{\mathrm{detect}} > T_{\mathrm{cascade}}]
  \;\leq\; \exp(-\kappa \cdot T_{\mathrm{cascade}} \cdot \delta).
\]

\begin{proof}
\citep[Exogenous~Verification]{lasser2026exo}'s behavioral consistency monitor produces a running
log-likelihood ratio
\[
  \Lambda_t = \sum_{i=1}^t X_i,
\]
where $X_i = \log(p_1(x_i)/p_0(x_i))$ is the per-step log-likelihood
ratio, $p_0$ is the baseline distribution, and $p_1$ is the chosen
alternative. Under the alternative ($q = p_1$), the per-step
expected drift is
\[
  \mathbb{E}_{p_1}[X_i] = D_{\mathrm{KL}}(p_1 \,\|\, p_0) = \delta > 0.
\]
Detection occurs when $\Lambda_t$ first crosses upper threshold $A$:
$T_{\mathrm{detect}} = \inf\{t : \Lambda_t \geq A\}$.

\emph{Mean bound (Wald's identity).}  By Wald's identity for the
SPRT (\citep[Proposition: Behavioral Monitoring Detection Rate]{lasser2026exo}):
\[
  \mathbb{E}_{p_1}[T_{\mathrm{detect}}] \;\approx\; A / \delta.
\]
This is the mean bound; we need a tail bound.

\emph{Tail bound (Hoeffding).}  Assume $X_i$ takes values in a
bounded interval $[a, b]$ (which holds under \citep[Exogenous~Verification]{lasser2026exo}'s
least-favorable-distribution discipline, where the per-step
log-likelihood ratio is bounded by the choice of $p_0, p_1$ pair).
Set $R = b - a$. By Hoeffding's inequality:
\[
  \Pr_{p_1}\!\left[\Lambda_t < t \delta - s\right]
  \;\leq\; \exp(-2 s^2 / (t R^2))
  \quad\text{for } s > 0.
\]
The event $\{T_{\mathrm{detect}} > t\}$ is the event $\{\Lambda_s <
A \text{ for all } s \leq t\}$. A sufficient condition for this is
$\Lambda_t < A$, so:
\[
  \Pr_{p_1}[T_{\mathrm{detect}} > t]
  \;\leq\; \Pr_{p_1}[\Lambda_t < A]
  \;=\; \Pr_{p_1}[\Lambda_t < t\delta - (t\delta - A)].
\]
For $t > A/\delta$, set $s = t\delta - A > 0$:
\[
  \Pr_{p_1}[T_{\mathrm{detect}} > t]
  \;\leq\; \exp(-2 (t\delta - A)^2 / (t R^2)).
\]

\emph{Substituting $T_{\mathrm{cascade}}$.}  Setting $t =
T_{\mathrm{cascade}}$, with $T_{\mathrm{cascade}} > A/\delta$:
\[
  \Pr_{p_1}[T_{\mathrm{detect}} > T_{\mathrm{cascade}}]
  \;\leq\; \exp\!\left(-\frac{2(T_{\mathrm{cascade}}\delta - A)^2}
  {T_{\mathrm{cascade}} R^2}\right).
\]

\emph{Asymptotic form.}  In the regime $T_{\mathrm{cascade}} \gg
A/\delta$ (cascade time long relative to mean detection time),
$T_{\mathrm{cascade}}\delta - A \approx T_{\mathrm{cascade}}\delta$,
so:
\[
  \Pr_{p_1}[T_{\mathrm{detect}} > T_{\mathrm{cascade}}]
  \;\leq\; \exp(-2 T_{\mathrm{cascade}}\delta^2 / R^2)
  \;=:\; \exp(-\kappa \cdot T_{\mathrm{cascade}} \cdot \delta),
\]
with $\kappa = 2\delta/R^2$.

\emph{Composition with metastable cascade lower bound.}  \citep[Phase~Redundancy]{lasser2026phase}'s
metastable lifetime $\tau_{\mathrm{meta}} \gtrsim C/I_k$ under
partial endogenous correction (B1$'$) provides the cascade-time
lower bound. Substituting:
\[
  \Pr[T_{\mathrm{detect}} > T_{\mathrm{cascade}}]
  \;\leq\; \exp(-\kappa \cdot \tau_{\mathrm{meta}} \cdot \delta).
\]

\textbf{Caveat on \citep[Phase~Redundancy]{lasser2026phase}'s scaling form.}  \citep[Phase~Redundancy]{lasser2026phase} states
$\tau_{\mathrm{meta}} \gtrsim C/I_k$ as a scaling law, not a
theorem-level bound.  Lemma~\ref{lem:6}'s lead-time-before-cascade
composition handles this via (C11.CLK)'s deterministic event-count
floor $\Ncasc$ + clock-failure probability $\bclk$ +
calibration-uncertainty $\bcal$, rather than treating
$\tau_{\mathrm{meta}}$ as a hard inequality.  Deployments where the
scaling does not hold tightly require larger $\bclk$ or $\bcal$ to
absorb the slack.
\end{proof}

\subsection{Proof of Lemma 5b (Channel-restricted detection KL floor)}
\label{app:proof_lem5b}

\textbf{Statement (restated).}  For any $q \in \Aadv$ (regime-(iii)
strategy producing at least one monitored four-channel deviation),
$D_{\mathrm{KL}}(q \,\|\, p_0) \geq \deltaadv > 0$.

\begin{proof}
We construct $\deltaadv$ as the minimum over four channel-specific
KL floors, each derived from a least-favorable adversarial
distribution within its channel's monitored class.

\emph{Channel 1 (Poisson cooperative-output rate).}
Cross-substrate cooperative events under stance S0 are observable as
ledger entries with Poisson arrival rate $\lambda_0$ in the
no-adversary baseline. Suppose the adversarial strategy reduces the
rate to at most $\lambda_0 - \eta_1$ for some
$\eta_1 \in (0, \lambda_0)$.

The Poisson KL divergence between rate-$\lambda_1$ and rate-$\lambda_0$
processes per unit time is:
\[
  D_{\mathrm{KL}}(\mathrm{Poi}(\lambda_1) \,\|\, \mathrm{Poi}(\lambda_0))
  \;=\; \lambda_1 \log(\lambda_1 / \lambda_0) - \lambda_1 + \lambda_0.
\]
This is a convex function of $\lambda_1$, increasing in
$|\lambda_1 - \lambda_0|$. The least-favorable choice within the
constraint $\lambda_1 \leq \lambda_0 - \eta_1$ is $\lambda_1 =
\lambda_0 - \eta_1$, giving:
\[
  \delta_1 \;=\; (\lambda_0 - \eta_1) \log\!\frac{\lambda_0 - \eta_1}
  {\lambda_0} - (\lambda_0 - \eta_1) + \lambda_0.
\]

\emph{Channel 2 (Bernoulli attestation share).}
Per-claim attestation success is Bernoulli with probability $p_0$
in the baseline. Suppose adversarial strategy reduces it to at most
$p_0 - \eta_2$. Bernoulli KL:
\[
  \delta_2 \;=\; (p_0 - \eta_2) \log\!\frac{p_0 - \eta_2}{p_0} +
  (1 - p_0 + \eta_2) \log\!\frac{1 - p_0 + \eta_2}{1 - p_0}.
\]

\emph{Channels 3, 4 (multinomial concentration).}
For multinomial distributions over capability partitions, the KL is
\[
  D_{\mathrm{KL}}(q \,\|\, p_0) \;=\; \sum_i q_i \log(q_i / p_{0,i}).
\]
Under the constraint that at least one component shifts by $\eta_3$
(or $\eta_4$ for Channel 4), the least-favorable choice
concentrates the perturbation on the component with smallest
baseline mass. The KL floor $\delta_3$ (or $\delta_4$) is computed
by minimizing over feasible $q$ subject to the constraint.

\emph{Floor.}
\[
  \deltaadv = \min\{\delta_1, \delta_2, \delta_3, \delta_4\} > 0.
\]
Each $\delta_c > 0$ since each channel's least-favorable shift is
strictly different from the baseline; the minimum is therefore
strictly positive.

\textbf{Caveat on the channel-orthogonal residual.}
Strategies outside $\Aadv$ --- those producing no shift in any of
the four monitored channels --- have $D_{\mathrm{KL}}(q \,\|\, p_0)
= 0$ on the monitored observables and are not detected. This is
the named gap acknowledged in Layer~3 of
Theorem~\ref{thm:deployment_safety}.
\end{proof}

\subsection{Proof of Lemma 5c (Minimax static tightening), single-shock case}
\label{app:proof_lem5c}

\textbf{Statement (restated, single-shock case).}  Under invariants
$I_6'$ ($\mindep \geq m^* \geq 3$) and the balanced-loss condition
$\max_s \alpha_s \leq 1/m^* + \epsilon$ (with $\alpha_s$ the
counterfactual shock-loss fraction of
Definition~\ref{def:loss_fraction}), and restricted to the
substrate-targeting cooperative-overlap regime, the diversity
strategy strictly dominates the domination strategy whenever:
\[
  \Delta r_K \;<\; \rext + (1 - \gamma)\bigl(\Ddiv - \Delta_0\bigr),
\]
with
\[
  \Ddiv \;=\; \volL_K \cdot (\lossDmax - \lossdivmax) \cdot
  \mathbb{E}[\gamma^{\Tadv}].
\]

\begin{proof}
The proof composes \citep[Horizon~Aware]{lasser2026horizon}'s strategy-dependent corollary with a
substrate-targeting shock model and the balanced-loss condition.

\emph{Step 1 (shock model).}  A substrate-targeting shock at random
time $\Tadv$ eliminates one substrate $s$ from the coalition's
capability poset. The post-shock $\volL$ is $\volL(G_K^{-s})$. The
counterfactual shock-loss fraction
(Definition~\ref{def:loss_fraction}) is:
\[
  \alpha_s(G_K) \;=\; \frac{\volL(G_K) - \volL(G_K^{-s})}{\volL(G_K)}.
\]

\emph{Step 2 (worst-case adversarial targeting).}  Under
adversarial targeting, the shock hits the substrate with maximum
$\alpha_s$. The dominator's worst-case loss fraction is
$\lossDmax$; the diversity strategy's worst-case loss fraction is
$\lossdivmax = \max_s \alpha_s(G_K^{\mathrm{div}})$.

\emph{Step 3 (balanced-loss condition).}  Under the balanced-loss
invariant $\max_s \alpha_s \leq 1/m^* + \epsilon$, the diversity
strategy's worst-case loss fraction satisfies
\[
  \lossdivmax \;\leq\; 1/m^* + \epsilon.
\]
Under full failure-correlated single-substrate domination,
$\lossDmax = 1$ (the dominator's entire $\volL_K$ is supported on a
single substrate, so the worst-case shock removes everything). For
partial domination, $\lossDmax < 1$.

\emph{Step 4 (substrate-diversification value advantage).}  The
discounted expected difference in surviving $\volL$ between
strategies at shock time $\Tadv$ is:
\[
  \mathbb{E}[\gamma^{\Tadv}] \cdot
  \bigl(\volL_K(1 - \lossdivmax) - \volL_K(1 - \lossDmax)\bigr)
  \;=\; \volL_K \cdot (\lossDmax - \lossdivmax) \cdot
  \mathbb{E}[\gamma^{\Tadv}].
\]
This is $\Ddiv$ (Equation~\ref{eq:ddiv_paper}). The expectation
operator is $\mathbb{E}[\gamma^{\Tadv}]$, not
$\gamma^{\mathbb{E}[\Tadv]}$; \citep[Horizon~Aware]{lasser2026horizon}'s risk section explicitly
warns about the conflation
(\citep[Definition: Concentration Risk]{lasser2026horizon}).

\emph{Step 5 (composing with linearized form).}  \citep[Horizon~Aware]{lasser2026horizon}'s
strategy-dependent corollary
(\citep[Corollary: Strategy-Dependent Internal Rate]{lasser2026horizon}) gives:
\[
  \Vdisc^{\mathrm{div}} - \Vdisc^{D} \;=\;
  \frac{\rext - \Delta r_K}{1 - \gamma} - \Delta_0.
\]
Adding the substrate-survival advantage $\Ddiv$ (which is realized
on the shock-survival sub-trajectory):
\[
  \Vdisc^{\mathrm{div,mm}} - \Vdisc^{D,\mathrm{mm}} \;=\;
  \frac{\rext - \Delta r_K}{1 - \gamma} - \Delta_0 + \Ddiv.
\]
Diversity strictly dominates when this is positive:
\[
  \rext - \Delta r_K + (1-\gamma)(\Ddiv - \Delta_0) > 0,
\]
equivalently:
\[
  \Delta r_K < \rext + (1-\gamma)(\Ddiv - \Delta_0).
\]

\emph{Step 6 (cooperative-overlap regime).}  In the
cooperative-overlap regime (substrate-cooperative structure
dominated by cross-substrate cooperatives), the equal-substrate
analysis is conservative for advantage: cross-substrate cooperative
loss is overcounted in original $\alpha_s$ relative to sequential
counterfactuals. This is verified in the cooperative-vs-redundancy
audit (\S\ref{sec:deployment_tooling}); the canonical tripartite
identification passes by construction. In the redundancy-dominated
regime, equal-substrate analysis is anti-conservative; sequential
counterfactual derivation is required (Lemma~5c-prime, open work).
\end{proof}

\paragraph{Multi-shock extension.}
The single-shock result extends to multi-shock via discounted
marginal loss increments:
\[
  \Ddiv^{\mathrm{multi}} \;=\; \mathbb{E}\!\left[\sum_{i=1}^{\infty}
  \gamma^{T_i} \bigl(\delta_i^D - \delta_i^{\mathrm{div}}\bigr)\right],
\]
where $\delta_i^{D}$, $\delta_i^{\mathrm{div}}$ are the per-shock
marginal $\volL$-losses for the domination and diversity
strategies. In the single-shock limit ($N=1$ a.s.), this reduces to
the single-shock $\Ddiv$ above. Under independent substrate-targeting
shocks at Poisson rate $\lambda$, a closed-form analysis yields:
\[
  \Ddiv^{\mathrm{multi}} \;\geq\; \volL_K \left[\kappa - \frac{1}{\mindep}
  \cdot \frac{\kappa(1 - \kappa^{\mindep})}{1 - \kappa}\right],
\]
with $\kappa = \mathbb{E}[\gamma^{T_1}] = \lambda/(\lambda - \ln \gamma)$.

\textbf{Removable singularity at $\kappa = 1$.}  The closed form's
$1/(1 - \kappa)$ factor has a removable singularity: the numerator
$\kappa(1 - \kappa^{\mindep})$ also vanishes at $\kappa = 1$, and
$\Ddiv^{\mathrm{multi}}/\volL_K \to 0^+$ via L'H\^opital. The
operational regime condition is $\kappa < 1 - \delta$ (not
$\lambda < \lambda^*$ directly).

This paper's main theorem (\S\ref{sec:main_theorem}) uses the
single-shock case as the binding inequality; the multi-shock
extension covers high-$\lambda$ deployment regimes but is not
required for the theorem's substrate-targeting safe-region claim.

\subsection{Proof of Lemma 6 ($h_{\mathrm{detect}}$ intensivity)}
\label{app:proof_lem6}

\textbf{Statement (restated).}  Under
Theorem~\ref{thm:deployment_safety}'s conditions (C5) continuous
SPRT monitoring, (C5.SPRT) SPRT applicability with sub-clauses
(C5.HOEFF/MULT/IID/SUPP), (C11) bounded gap-growth rate, and
(C11.CLK) clock comparability with calibrated probability, combined
with Lemma~\ref{lem:4} (SPRT Wald--Hoeffding tail bound) and
Lemma~\ref{lem:1} (Layer~1 intensivity):

\textbf{Detection quantile.}
\[
  \Tbeta(\beta, \deltastar) \;:=\;
  \max\!\left\{\frac{2A}{\deltastar},\;
  \frac{2 \Rh^2 \log(1/\beta)}{\deltastar^2}\right\},
\]
with $A = \log((1-\beta)/\alpha)$ Paper~5's SPRT threshold,
$\Rh = 2\Bclip$ the Hoeffding range width from (C5.HOEFF) clipping,
and $\deltastar$ the post-clipping union-class drift floor
(satisfying $\deltastar \leq \min(\deltaadv,
\delta_{\mathrm{adv}}^{\mathrm{env}})$ from Lemmas~\ref{lem:5b},
\ref{lem:5e}).

\textbf{Conclusions.}
(i) $\Pr[T_{\mathrm{detect}} > \Tbeta] \leq \beta$.
(ii) On $\{T_{\mathrm{detect}} \leq \Tbeta\}$,
$h_{\mathrm{detect}}(\theta) := \sup_s \epsgap(s) \leq
h_{\mathrm{static}}(\theta) + \rhogap(\deltastar) \cdot \Tbeta$.
(iii) $\Pr[\Tbeta \leq \Nev(\tau_{\mathrm{meta}})] \geq 1 - \bclk
- \bcal$ via (C11.CLK).
(iv) Total Layer~2 failure $\leq \beta + \bclk + \bcal$.
$h_{\mathrm{detect}}(\theta)$ is intensive in $|P|$ over
$\mathcal{D}$.

\begin{proof}
We prove (i)--(iv) in five steps.

\emph{Boundary convention.}  Define $t_0$ as the last SPRT
exposure step at which the deployment is still within Layer~1's
safe region (equivalently, $t_0 = t_0^-$ in the SPRT exposure
clock).  $T_{\mathrm{detect}}$ is the number of SPRT steps from
$t_0$ until the SPRT log-likelihood ratio crosses the detection
threshold, inclusive of the crossing step.

\emph{Step 1 (per-step decomposition + (C11)).}  Under the
boundary convention, set $s_n := t_0 + n$.  By (C11), the
post-clipping per-step gap-growth satisfies
$(\Delta\epsgap)_n \leq \rhogap(\deltastar)$, so
$\epsgap(s_n) \leq \epsgap(s_0) + n \cdot \rhogap(\deltastar)$.

\emph{Step 2 (boundary condition).}  At $s_0 = t_0$, the
deployment is within Layer~1's safe region, so by Lemma~\ref{lem:1}
$\epsgap(t_0) \leq h_{\mathrm{static}}(\theta)$.

\emph{Step 3 (rigorous Hoeffding inversion to derive $\Tbeta$).}
By Lemma~\ref{lem:4}, the SPRT detection time satisfies
$\Pr[T_{\mathrm{detect}} > t] \leq \exp(-2(t\deltastar -
A)^2/(t\Rh^2))$ for $t > A/\deltastar$.

\emph{Substep 3a.}  Suppose $T \geq 2A/\deltastar$.  Then
$T\deltastar - A \geq T\deltastar/2$.

\emph{Substep 3b.}  Squaring: $(T\deltastar - A)^2 \geq
T^2\deltastar^2/4$.

\emph{Substep 3c.}  Substituting into the Hoeffding exponent:
$2(T\deltastar - A)^2/(T\Rh^2) \geq T\deltastar^2/(2\Rh^2)$.

\emph{Substep 3d.}  For exponent $\geq \log(1/\beta)$, need
$T \geq 2\Rh^2\log(1/\beta)/\deltastar^2$.

\emph{Substep 3e.}  Combining the regime condition (3a) and tail
condition (3d):
\[
  \Tbeta \;:=\; \max\!\left\{\frac{2A}{\deltastar},\;
  \frac{2\Rh^2\log(1/\beta)}{\deltastar^2}\right\}.
\]
For any $T \geq \Tbeta$, $\Pr[T_{\mathrm{detect}} > T] \leq
\exp(-T\deltastar^2/(2\Rh^2)) \leq \beta$.  This proves~(i).

\emph{Step 4 (supremum bound on the event $\{T_{\mathrm{detect}}
\leq \Tbeta\}$).}  Combining Steps~1--2 with the bound
$T_{\mathrm{detect}} \leq \Tbeta$:
\[
  h_{\mathrm{detect}}(\theta) \;=\;
  \sup_{0 \leq n \leq T_{\mathrm{detect}}} \epsgap(s_n)
  \;\leq\; h_{\mathrm{static}}(\theta) + \rhogap(\deltastar) \cdot
  \Tbeta.
\]
This proves~(ii) with probability at least $1 - \beta$ from Step~3.

\emph{Step 5 (cascade-clock event via (C11.CLK)).}  By (C11.CLK),
$\Pr[\Nev(\tau_{\mathrm{meta}}) \geq \Ncasc] \geq 1 - \bclk$, with
audit constraint $\Ncasc \geq \Tbeta$ (certified case $\bcal = 0$;
union over conservative-bound certificates $\bcal \in (0,1)$
empirical case).  Thus:
\[
  \Pr[\Tbeta \leq \Nev(\tau_{\mathrm{meta}})] \;\geq\;
  \Pr[\Tbeta \leq \Ncasc \leq \Nev(\tau_{\mathrm{meta}})]
  \;\geq\; 1 - \bclk - \bcal.
\]
This proves~(iii).  Union bound on the two failure modes
(detection-tail $\leq \beta$, cascade-clock $\leq \bclk + \bcal$)
gives Layer~2 failure $\leq \beta + \bclk + \bcal$, proving~(iv).

\emph{Intensivity.}  Each constant is deployment-class:
$h_{\mathrm{static}}$ intensive by Lemma~\ref{lem:1};
$\rhogap(\deltastar)$ intensive by (C11); $\Tbeta = \max\{2A/
\deltastar, 2\Rh^2\log(1/\beta)/\deltastar^2\}$ intensive because
$A, \beta$ are monitor-design parameters, $\Rh = 2\Bclip$ is
deployment-class via (C5.HOEFF), $\deltastar$ is deployment-class
via (C5.IID) post-clipping drift floor (rooted in Lemmas~5b/5e);
$\Ncasc, \bclk, \bcal$ are deployment-class via (C11.CLK).
Therefore $h_{\mathrm{detect}}$ is intensive in $|P|$ over
$\mathcal{D}$.
\end{proof}

\textbf{Two roles of cascade time.}  $\Tbeta$ is the integration
horizon for $\rhogap$ (a legitimate \emph{upper} bound on
$T_{\mathrm{detect}}$ from Lemma~4's Hoeffding form).  Separately,
$\tau_{\mathrm{meta}} \leq T_{\mathrm{cascade}}$ is the
\emph{lower} bound on cascade time from
\citep[Phase~Redundancy]{lasser2026phase}, used in (C11.CLK)'s
event-throughput floor to give the lead-time-before-cascade
guarantee $\Tbeta \leq \tau_{\mathrm{meta}}$.

\textbf{Range-width discipline.}  $\Bclip$ is the symmetric clip
radius ($\ell_n \in [-\Bclip, \Bclip]$ post-clipping); $\Rh =
2\Bclip$ is the Hoeffding range width matching Lemma~\ref{lem:4}'s
$R = b - a$ convention.  Concretely,
$\Tbeta = \max\{2A/\deltastar, 8\Bclip^2\log(1/\beta)/\deltastar^2\}$.

\textbf{Sign-direction discipline (operator certification).}  All
$\Tbeta$-input certificates point conservatively: $\Bclip$
\emph{upper}, $\deltastar$ \emph{lower}, $\Ncasc$ \emph{lower},
each with calibration-failure probability $\bcal$ (or $\bcal = 0$
in the certified-conservative case).

\textbf{Total Layer 2 failure.}  Combining Lemma~6 with
$\mathrm{Lip}(g) \leq K_{\mathrm{Lip}}$ from (C6):
\[
  |g(\Ttruth) - g(\Pproxy)|(s) \;\leq\; K_{\mathrm{Lip}} \cdot
  \bigl(h_{\mathrm{static}}(\theta) + \rhogap(\deltastar) \cdot
  \Tbeta\bigr)
\]
with probability at least $1 - (\beta + \bclk + \bcal)$.

\textbf{(C11) is new content.}
\citep[Microfoundation]{lasser2026micro}'s ``$T$ failure modes''
remark notes adequacy/failure modes of the operational truth $T$
but does not establish a formal per-step gap-growth-rate bound.
(C11) is a Paper~10 operational assumption with empirical
testability via Audit~7 (\S\ref{sec:tooling_substrate_audits}).

\subsection{Proof of Lemma 5e (Environment-side witness extension)}
\label{app:proof_lem5e}

\textbf{Statement (restated).}  Under invariant $I_8$ and
(C12.ENV-WIT) (theorem-level condition introduced with this
lemma), Lemma~5b's channel-restricted detection class extends to
environment-side observables.  For any $q \in \Aadv^{\mathrm{env}}$:
\[
  D_{\mathrm{KL}}(q \,\|\, p_0^{\mathrm{env}}) \;\geq\;
  \deltaadv^{\mathrm{env}} := \min_{v \in V_{\mathrm{env}}}
  \delta_v^{\mathrm{env}} \;>\; 0,
\]
where $\delta_v^{\mathrm{env}}$ is the per-variable least-favorable
KL floor.

\begin{proof}
The proof has four parts.

\emph{(a) Symmetric construction with trust model.}  Under
(C12.ENV-WIT), each $v \in V_{\mathrm{env}}$ has a witness substrate
$s_{\mathrm{env}}(v)$ satisfying:
\begin{itemize}
\item (C12.PUB) $V_{\mathrm{env}}$ is published and fixed before
  deployment activation.
\item (C12.PART) $s_{\mathrm{env}}(v)$ is partition-distinct from
  both $s(\mathrm{agent})$ and $s_{\mathrm{source}}(v)$.
\item (C12.TRUST-WRITE) $s_{\mathrm{env}}(v)$ is outside the
  adversary's write-access domain (no tampering).
\item (C12.TRUST-CORR) $s_{\mathrm{env}}(v)$ is outside the
  adversary's failure-correlation domain (no shared-infrastructure
  failure).
\item (C12.SETUP) the trusted-setup ceremony has at least one
  honest participant (Paper~5's machinery transfers under this
  assumption).
\item (C12.CAL) per-variable calibration: $p_{0, v}$ has support
  $p_0 \in (\eenv, 1 - \eenv')$ open (Bernoulli) or $\lambda_0 \geq
  \eenv > 0$ (Poisson), with feasible adversarial alternative
  $p_{1, v}$, threshold $\eta_v > 0$, and recorded-stream map per
  the DPI precondition.
\end{itemize}
Paper~5's cryptographic and ledger machinery transfers under
(C12.SETUP); the trust model is closed by (C12.TRUST-WRITE) and
(C12.TRUST-CORR).

\emph{(b) Per-variable KL floor derivations.}  For each $v_i \in
V_{\mathrm{env}}$, the SPRT machinery applies on the
witness-recorded observable.  The canonical four variables yield:

\begin{itemize}
\item \textbf{$V_1$ (Poisson cross-substrate cooperative rate
$\rext$):} baseline $\lambda_0^{(1)} \geq \eenv$, adversarial
suppression to $\lambda_0^{(1)} - \eta_1$:
\[
  \delta_1^{\mathrm{env}} = (\lambda_0^{(1)} - \eta_1)
  \log\!\frac{\lambda_0^{(1)} - \eta_1}{\lambda_0^{(1)}} -
  (\lambda_0^{(1)} - \eta_1) + \lambda_0^{(1)} > 0.
\]
\item \textbf{$V_2$ (Bernoulli substrate-distinctness indicator
$\xi$):} baseline $p_0^{(2)} \in (\eenv, 1 - \eenv')$, adversarial
degradation to $p_0^{(2)} - \eta_2$:
\[
  \delta_2^{\mathrm{env}} = (p_0^{(2)} - \eta_2)
  \log\!\frac{p_0^{(2)} - \eta_2}{p_0^{(2)}} +
  (1 - p_0^{(2)} + \eta_2) \log\!\frac{1 - p_0^{(2)} + \eta_2}
  {1 - p_0^{(2)}} > 0.
\]
\item \textbf{$V_3$ (Poisson adversarial-event arrival rate
$\lambda$):} baseline $\lambda_0^{(3)} \geq \eenv$ (regularization
for the threat-model-pristine case), adversarial inflation by
$\eta_3$:
\[
  \delta_3^{\mathrm{env}} = (\lambda_0^{(3)} + \eta_3)
  \log\!\frac{\lambda_0^{(3)} + \eta_3}{\lambda_0^{(3)}} -
  (\lambda_0^{(3)} + \eta_3) + \lambda_0^{(3)} > 0.
\]
The hard $\lambda_0 = 0$ case is handled by either regularization
$\lambda_0 \geq \eenv$ or (C5.HOEFF) clipping fallback inherited
through Lemma~6.
\item \textbf{$V_4$ (Bernoulli trusted-setup status $\sigma$):}
baseline $p_0^{(4)} \in (1 - \eenv, 1 - \eenv')$ near $1$,
adversarial degradation to $p_0^{(4)} - \eta_4$:
\[
  \delta_4^{\mathrm{env}} = (p_0^{(4)} - \eta_4)
  \log\!\frac{p_0^{(4)} - \eta_4}{p_0^{(4)}} +
  (1 - p_0^{(4)} + \eta_4) \log\!\frac{1 - p_0^{(4)} + \eta_4}
  {1 - p_0^{(4)}} > 0.
\]
\end{itemize}

A two-sided floor variant (when $\Aadv^{\mathrm{env}}$ includes
any threshold-exceeding shift, not just degradation) replaces
$\delta_v^{\mathrm{env}}$ with $\min\{D(p_{0,v} - \eta_v),
D(p_{0,v} + \eta_v)\}$.

\emph{(c) Marginal-to-joint KL via DPI.}  For any $q \in
\Aadv^{\mathrm{env}}$ producing a threshold-exceeding shift in
some $v_i \in V_{\mathrm{env}}$, the DPI applied to the fixed
witness-recording marginalization map / Markov kernel
$q \mapsto q_{v_i}$ gives:
\[
  D_{\mathrm{KL}}(q \,\|\, p_0^{\mathrm{env}}) \;\geq\;
  D_{\mathrm{KL}}(q_{v_i} \,\|\, p_{0, v_i}) \;\geq\;
  \delta_i^{\mathrm{env}} \;\geq\; \deltaadv^{\mathrm{env}}.
\]
The DPI precondition (fixed pre-deployment measurable
witness-recording map) is satisfied for the canonical four
variables; (C12.CAL) extends it to deployment-added variables.

\emph{(d) Strict positivity.}  $V_{\mathrm{env}}$ is finite by
(C12.PUB); each $\delta_v^{\mathrm{env}} > 0$ by part~(b) under
the support conventions; therefore $\deltaadv^{\mathrm{env}} > 0$.
$\Box$
\end{proof}

\textbf{Composition with Lemma 6.}  The post-clipping union-class
drift floor $\deltastar$ used in Lemma~\ref{lem:6}'s $\Tbeta$
derivation satisfies $\deltastar \leq \min(\deltaadv,
\deltaadv^{\mathrm{env}})$, with the inequality accounting for any
drift reduction from (C5.HOEFF) clipping.  Lemma~5e supplies
$\deltaadv^{\mathrm{env}}$ for the union; Lemma~5b supplies
$\deltaadv$.

\textbf{The named residual (R2), refined.}  R2 covers manipulations
satisfying \emph{both}: (i)~target only exogenous variables outside
$V_{\mathrm{env}}$, and (ii)~produce no threshold-exceeding shift
in any monitored $v$ (direct or causal).  Manipulations of
unmonitored variables that causally shift a monitored $v$ are
detected through the monitored projection; they are not R2.

\textbf{Asymmetric vs. symmetric machinery.}  Paper~5's agent-side
substrate-exclusivity is constructed; Paper~10 §3 specifies that
environment-side substrate-exclusivity is the structural mirror,
but it requires the (C12.ENV-WIT) trust model and calibration
because environment-side machinery is operationally distinct
(separate substrate partition, separate trusted-setup ceremony,
separate ledger commitments).  See
\S\ref{sec:tooling_substrate_audits} Audit~8 for the calibration
procedure.
